\section{秦昭襄王：伊阙之后，如何把胜利变成次序}

秦武王猝卒时，王位落到年轻的异母弟手中。前307年的\eventanchor{WQ-0307-QIN-SUCCESSION}{秦·昭襄王即位}，并不只是年表上一笔继承：宣太后和魏冉所连结的宫廷、外戚与将领网络，先替新王稳住了朝廷，也使军政大权在相当长的时间里不能只用“王命”二字来概括。昭襄王的漫长在位由此开始。他所继承的，是商鞅以来能够征发户口、粮食和兵员的秦国；他所必须学会的，则是如何在太后、相国、宗室和战将各有资源的格局里，使这些资源服从一条较持久的东进路线。

伊水与洛水会合处的峡口，正是这种路线的一次试验场。前293年的\eventref{WQ-0293-YIQUE}{战国·伊阙之战}中，白起所部击溃韩、魏联军，韩魏面向洛阳和函谷关的屏障因而更薄。传世史书把斩获写得极多，又把战果归于白起一人的临阵才能；可以较稳妥把握的是秦军在这一通道取得决定性优势，至于杀伤数字、战斗展开的日次和各军具体位置，仍主要来自后出的叙事，不宜当作可逐项核验的战报。峡口、城邑和渡口比华丽的战功更能说明此战的分量：它让秦更容易把关中仓廪和兵员送到东面，也迫使韩、魏在一条越来越狭窄的防线上安排援军。

\begin{center}
\begin{tikzpicture}[x=.9cm,y=.9cm,font=\small]
  \fill[paper] (0,0) rectangle (10.8,5.8);
  \draw[gray!55,line width=1pt] (3.6,.5) .. controls (4.3,1.7) and (3.7,3.7) .. (4.5,5.4);
  \node[rotate=76,text=gray!70] at (3.95,4.1) {崤函山地};
  \draw[blue!55,line width=1pt] (.4,4.6) .. controls (3.2,4.25) and (6.7,4.95) .. (10.2,4.45);
  \node[text=blue!70!black,above] at (6.7,4.65) {黄河方向};
  \draw[blue!45,line width=.8pt] (4.6,.5) .. controls (5.0,1.4) and (5.6,2.35) .. (5.5,3.65);
  \node[text=blue!70!black,right] at (5.5,1.35) {伊水、洛水方向};

  \fill[dynasty!24] (.8,1.2) -- (3.25,1.0) -- (3.45,4.55) -- (1.25,4.9) -- cycle;
  \node[align=center] at (2.05,2.95) {秦\\关中};
  \fill[yellow!24] (4.15,2.75) -- (5.15,2.6) -- (5.35,3.75) -- (4.35,3.95) -- cycle;
  \node at (4.75,3.28) {韩};
  \fill[green!20] (5.85,3.0) -- (8.25,2.9) -- (8.45,4.55) -- (6.05,4.7) -- cycle;
  \node at (7.15,3.78) {魏};

  \fill[timeline] (3.85,2.95) circle (.12);
  \node[below,align=center] at (3.85,2.88) {函谷\\方向};
  \fill[timeline] (5.25,2.05) circle (.14);
  \node[below,align=center] at (5.25,1.96) {伊阙\\前293年};
  \fill[source] (5.62,3.72) circle (.1);
  \node[above,align=center] at (5.62,3.74) {洛阳\\方向};

  \draw[-{Stealth[length=2mm]},ultra thick,color=dynasty] (2.9,2.65) -- (5.05,2.1);
  \draw[-{Stealth[length=2mm]},thick,color=yellow!60!black] (4.68,2.75) -- (5.12,2.2);
  \draw[-{Stealth[length=2mm]},thick,color=green!60!black] (6.22,3.05) -- (5.4,2.2);
  \node[align=center,text=dynasty] at (3.55,1.92) {秦军东出\\与补给方向};
  \node[align=center,text=source] at (7.85,2.25) {韩、魏联军\\向峡口集结};
\end{tikzpicture}

\smallskip
\small\textit{图\quad 伊阙—函谷走廊示意：据《史记·秦本纪》《白起王翦列传》及《资治通鉴》呈现前293年伊阙之战所处的关中东出口、伊洛峡口与韩魏防线之间的关系。山地、河流和地点用于说明通道，不复原军队日次部署或补给线；各政治体边界均为约略示意。}
\end{center}


但伊阙没有使秦从此只须向东推行。齐、楚仍足以牵制，韩、魏、赵的临时合从也会令函谷关以外的胜利变得昂贵。魏冉经营的封邑、宾客和军功关系，曾为这种扩张提供可调度的人与物；同一条件也留下王权如何约束相权的难题。结构性原因不在于秦廷忽然有了一位能看穿天下的谋士，而在于连续战争已经使每一次出兵都同时成为对粮道、盟约和宫廷指挥权的考验。

范雎进入秦廷，正落在这道考题最紧的时候。关于他的魏国遭遇、改名入秦以及荐引经过，《史记》所写极具人物故事的张力，《战国策》的说辞又有不同层次；这些材料能够说明一个外来游士如何被安置进秦国权力中心，不能让后人把每一段受辱、密谈和报复复原为现场。前266年前后的\eventanchor{WQ-0266-FANJU}{秦·范雎入秦与远交近攻}，较可靠的核心是昭襄王借此调整权力，魏冉终于失势，范雎居相；人事转折所改变的，不只是一个职位，而是宫廷中谁能够把邻国的远近、战争的成本和王的名义连成一项政策。

\begin{textwindow}[《史记》中的一句方略]
“王不如远交而近攻，得寸则王之寸也，得尺亦王之尺也。”

这句见于\sourcebook{史记·范雎蔡泽列传}，后来成为解释秦国外交的钥匙。它以排比和断语组织成一场成功游说，很可能已经受司马迁叙事的剪裁；《战国策》同类说辞亦不宜逐字视为秦廷记录。可是，它准确点出了可由战争地理检验的选择：远方大国即使暂时修好，也未必能替秦守住新得城邑；韩、魏、赵等近邻的土地一旦夺取，才较可能同既有县邑、道路和征发网络相连。“远交近攻”因此不是包治天下的口号，而是把外交的暂时让步同近处的持续吞并配合起来的计算。
\end{textwindow}

这套计算并未使赵自动退让。赵在太行以东的城邑、北方扩张所得的人口以及通向邯郸的道路，使它成为秦在中原北缘最难绕开的对手；秦军的推进也一再遇到抵抗。到前262年，秦攻取野王，韩国上党同本国腹地的交通被截断。上党郡守不愿降秦而求援于赵，赵廷接收上党，于是\eventanchor{WQ-0262-SHANGDANG}{战国·上党归赵}把原先的韩秦战事转成秦赵之间的正面争夺。这里没有一方被动“卷入”：韩国地方官在降与求援之间作出选择，赵要权衡得到山地城邑和承受秦军压力的代价，秦则不能容忍东出通道旁出现一个由赵接手的敌对据点。

直接后果是前线沿太行山地延长，双方都得把粮秣、役夫、守城兵和朝廷的耐心运往山口。中期影响才是\eventref{WQ-0260-CHANGPING}{战国·长平之战}的形成：它不是范雎一句话或白起一场胜利必然推出的终局，而是近攻的战略偏好、上党归属的触发因素、秦赵各自可动员的资源与盟国难以及时协同共同叠加的结果。

\subsection{在昭襄王之前：把西南与函谷关接进同一张网}

昭襄王继承的疆域并非只靠函谷关积累。前316年秦军入蜀，石牛道等山路、栈道和谷口限制大军行速，\eventanchor{WQ-0316-QIN-STONE-CATTLE-ROAD}{战国·秦军经蜀道入蜀}；“石牛”传说的细部不宜当作行军日记，但山地转运的困难确属这场征服的条件。蜀相陈壮在秦军压境时杀蜀王，\eventanchor{WQ-0316-CHENZHUANG-KILLS-SHU}{战国·陈壮杀蜀王}，内部政变加速了抵抗的崩解；巴王降秦，\eventanchor{WQ-0316-BASHU-SURRENDER}{战国·巴王降秦}。降附如何协商、地方首领保留多少统治并无充分材料，却可确认巴蜀的山川、人口和峡江道路后来成为秦须经营、也能调用的西南资源。约前300年，秦在谷口、亭障与转输点反复维护蜀道，\eventanchor{WQ-0300-QIN-SHU-ROAD-MILESTONES}{战国·秦经营蜀道转输}；节点名称和修筑年次并不完备，只能谨慎说关中和成都平原开始能够进行季节性的物资流动。

前311年惠文王卒、武王即位，\eventanchor{WQ-0311-QIN-WUWANG}{战国·秦武王即位}。新君前309年置左右丞相，\eventanchor{WQ-0309-QIN-CHANCELLORS}{战国·秦武王初置左右丞相}，并不等于已经出现后世意义的宰相制，却显示宗室、客卿和军队需要被分途协调。甘茂远攻宜阳前以息壤之盟要求武王勿疑，\eventanchor{WQ-0308-QIN-XIRANG-PLEDGE}{战国·甘茂以息壤之盟求信}；盟誓的完整辞令来自列传，不是现场契约，仍能看出远征的成败也取决于君将之间能否保持足够久的信任。宜阳终于被取，\eventanchor{WQ-0308-YIYANG}{战国·秦取宜阳}，将战果、守备和韩地通道一并推到更靠东的位置。

武王猝卒后，王位转到昭襄王，\eventref{WQ-0307-QIN-SUCCESSION}{秦·昭襄王即位}。前306年诛庶长壮等政变相关者，\eventanchor{WQ-0306-QIN-COURT-PURGE}{战国·昭襄王清理政变相关者}，涉案范围无从逐人核对；较可靠的是摄政结构得以稳住。宣太后、魏冉的网络在短期内提供了军政连续性，也使年轻的王必须在外戚、宗室、客卿和战将之间取得真正的任免权。甘茂以“曾参杀人”劝武王坚持攻宜阳的说辞，\eventanchor{WQ-0308-GANMAO}{战国·甘茂以曾参杀人之喻求信}，也应看作列传以故事解释君将互信的写法，而非一段可逐字核验的军令。

\subsection{攻城之后，如何把前沿留住}

白起的上升正发生在这种朝廷结构中。前297年孟尝君失秦相位，秦廷重新限制齐客卿介入军政，\eventanchor{WQ-0297-QIN-MENGCHANG-DISMISSAL}{战国·孟尝君失秦相与客卿风险}；逃离与失相的细节多有传奇加工，齐秦关系转为公开对抗则是可见的后果。前294年前后，白起以对魏攻城进入主将序列，\eventanchor{WQ-0294-BAIQI-APPOINTMENT}{战国·白起进入主将序列}，并取曲沃、邢丘等地，\eventanchor{WQ-0294-BAIQI-QUWO}{战国·白起取曲沃邢丘}。城名归属与最初官号均有考辨，不能把军功履历补成完整档案；秦军能够反复向魏的城邑网络施压，才是更可靠的事实。前292年取宛后，秦把南阳北部作为持续进攻的据点，\eventanchor{WQ-0292-QIN-WAN-COMMANDERY}{战国·秦取宛作为进攻据点}；同年取垣又旋即归还，\eventanchor{WQ-0292-WEI-YUAN-RETURN}{战国·秦取垣后归还}。后一举的动机不必强说为某种高明计策，却说明占领、还城与外交威慑可以交替使用，使魏难以据一城判断秦的下一步。

司马错取轵、邓后的山口与河道压力，\eventanchor{WQ-0291-QIN-SHANGDANG-CORRIDOR}{战国·秦压迫上党通道}，以及沿此方向勘察粮运条件的谨慎推断，\eventanchor{WQ-0291-QIN-HAN-ROUTE-SURVEY}{战国·秦军勘察韩地进攻轴线}，都不该化作一份失传的测绘图。它们说明前线不靠临阵选择维系：秦须在每次攻城前判断何处山口能接上后方，韩魏则须在同一片狭窄地带分兵。前298年齐韩魏攻秦时，秦依函谷关、崤函谷道与关中仓储延缓联军，\eventanchor{WQ-0298-QIN-HANGU-DEFENSE}{战国·秦依函谷关防御三国联军}；守将、兵力与关塞部署未有完整记录，不能把地形写成自动取胜的屏障。前289年司马错在河道、桥梁与城邑间行动，\eventanchor{WQ-0289-QIN-HEXI-RIVERWORK}{战国·司马错经营河道攻城}，又经营两岸渡口，\eventanchor{WQ-0289-QIN-RIVER-CROSSINGS}{战国·秦经营黄河渡口}。桥梁、决河和城址细节各有异说，河防被切开而秦能多向机动的后果，却比传奇战法更重要。

前287年昭襄王巡汉中、上郡、北河，\eventanchor{WQ-0287-QIN-NORTHWEST-TOUR}{战国·昭襄王巡行西北}；这次\eventanchor{WQ-0287-QIN-NORTHERN-TOUR}{战国·昭襄王察看汉中上郡北河}巡行的路线和具体举措不可尽定，亦不应想象为一张由王亲自批示的边区图。同一时期，秦廷透过义渠王的往来与联姻观察北西边情，\eventanchor{WQ-0287-YIQU-INTELLIGENCE}{战国·秦廷经营义渠信息}。这种后出的宫廷叙事不能让人推断私密关系的动机，却提示北方边患已进入王廷的战略盘算。前285年河东分置九县，\eventanchor{WQ-0285-QIN-HEDONG-NINE-COUNTIES}{战国·秦河东置九县}，即\eventanchor{WQ-0285-QIN-HEDONG-COUNTIES}{战国·河东以县治固化占领}，又以户籍、田界安排新占地的赋役守备，\eventanchor{WQ-0285-QIN-HEDONG-HOUSEHOLDS}{战国·秦河东九县编户守备}；县名、边界和户数不明，前沿从战场转为常设行政支点的方向却十分清楚。

\subsection{北方吞并与南方压力}

义渠的处置使北方边线从结交转向兼并。前272年秦诱杀义渠王并吞其地，\eventanchor{WQ-0272-YIQU}{战国·秦吞义渠}；宣太后以婚姻安抚、后在甘泉诱杀的故事，\eventanchor{WQ-0272-XUANTAIHOU-YIQU-ENTRAPMENT}{战国·甘泉诱杀义渠王}，来自后出史书，不能据此补写个人情感或预谋的全部过程。较可确认的是首领被除之后，秦在北地、陇西安排郡县、戍守和移民，\eventanchor{WQ-0272-YIQU-COMMANDERIES}{战国·义渠地的郡县戍守}。新获得的纵深也意味着管理牧地、交通和新附部众的长期成本，而非一块自动服从的北方。

南面则是另一套地理。前281年前后秦在汉水上游布置据点，\eventanchor{WQ-0281-QIN-CHU-HANSHUI}{战国·秦沿汉水压迫楚防线}；前280年取黔中后在河谷山地部署守军，\eventanchor{WQ-0280-QIN-QIANZHONG-MILITARY-SETTLEMENT}{战国·秦在黔中部署守备}。驻军规模、县治与控制深度都难确定，不能把后世郡界移回此时。前277年取巫、黔中时依峡江与山道转输，\eventanchor{WQ-0277-QIN-WUSHAN-LOGISTICS}{战国·秦军巫黔补给}，又利用水路把战果向楚西南延展，\eventanchor{WQ-0277-QIN-WUSHAN-RIVER-ROUTE}{战国·秦利用峡江水路}。这些补给线正说明，楚的纵深虽大，亦须在汉水、峡江和郢都间分配防守。

白起攻郢后受封武安君，\eventanchor{WQ-0278-BAIQI-WUAN}{战国·白起封武安君}，其军功与朝廷影响同时上升，\eventanchor{WQ-0278-BAIQI-MERIT-REWARD}{战国·白起破郢受赏}。封赏仪式和封邑不可实指，白起的声望能够在秦廷政策中占据重量却无须怀疑。秦又向韩国刚寿等南阳城邑施压，\eventanchor{WQ-0274-QIN-HAN-GANGSHOU}{战国·秦压迫韩国刚寿}；前264年前后攻韩、拔陉城等要点，\eventanchor{WQ-0264-BAIQI-HAN}{战国·白起攻韩拔陉}。城名与年份有异文，但持续压力使韩国不得不把有限兵粮用于西线。

前273年秦将上庸等地合为郡、迁南阳免臣，\eventanchor{WQ-0273-QIN-SHANGYONG-COMMANDERY}{战国·秦置上庸等郡}；前272年又初置南阳郡，\eventanchor{WQ-0272-QIN-NANYANG-COMMANDERY}{战国·秦置南阳郡}，并尝试整合汉水北部粮道与新占居民，\eventanchor{WQ-0272-QIN-NANYANG-GRAIN}{战国·南阳粮道与新占居民}。早期辖境和“免臣”身份仍待考，重要的是行政尝试同对韩、控楚的后勤需要同步出现，而不是先有一套完善郡制再拿去征服。

\subsection{外戚退场，王权重新选择进攻次序}

范雎入秦后的权力转折，并非只是一篇游说辞的胜利。约前268年魏冉出居陶，\eventanchor{WQ-0268-WEIRAN-EXILE}{战国·魏冉出陶}，宣太后的影响亦收缩，\eventanchor{WQ-0268-XUANTAIHOU-RETIREMENT}{战国·宣太后政治影响收缩}；罢相步骤与个人动机在《史记》自身也有先后差异，不能写作一次毫无阻滞的宫廷清洗。更稳妥的是，范雎居相后重组王权的决策渠道，\eventanchor{WQ-0268-QIN-FANJU-COURT}{战国·范雎重组秦廷决策}，并把先压韩赵魏的战略框架反复表述为远交近攻，\eventanchor{WQ-0268-QIN-FANJU-STRATEGY-FRAME}{战国·范雎的近攻战略框架}。这种说法不替代每次具体决策，却使近邻的城邑、道路和粮运成为王廷更优先计算的对象。

宫廷的继承链也在同一时期变动。前267年太子悼在魏卒、归葬芷阳，\eventanchor{WQ-0267-QIN-DAOTAI-RETURN}{战国·太子悼卒归芷阳}，秦魏之间的质子安排受冲击，\eventanchor{WQ-0267-QIN-HEIR-HOSTAGE-NETWORK}{战国·太子悼之死与质子网络}。死因、居魏目的和家族排行无从确证；秦廷必须重新审视安国君诸子与外戚的影响，\eventanchor{WQ-0267-AN-GUO-SUCCESSION-REVIEW}{战国·秦廷重审安国君继承}。前265年立安国君为太子、宣太后卒、魏冉出陶，\eventanchor{WQ-0265-QIN-ANGUO-CROWN-PRINCE}{战国·安国君立太子}；华阳夫人与异人周边的联盟开始显形，\eventanchor{WQ-0265-QIN-ANGUO-ALLIANCE}{战国·安国君继承联盟形成}。后来的吕不韦故事多有传记加工，仍不妨碍人们看见：统一前的王位政治，也由质任、婚姻与外戚关系预先塑形。

秦在魏地新占城邑还须面对旧乡里、守卒和官吏，\eventanchor{WQ-0266-QIN-WEI-LOCAL-ELITES}{战国·秦接管魏地地方关系}。接管方式没有文书可复原，不能凭后来的秦简替这一年填表；恰恰因为这类摩擦无法被一个胜利数字遮住，近攻才是一项需要持续投入行政的策略。

\subsection{长平不是一日决战}

前263年至前262年秦取野王，截断韩国与上党交通，\eventanchor{WQ-0263-YEWANG}{战国·秦取野王}；五大夫贲又取韩十城，\eventanchor{WQ-0262-QIN-HAN-TEN-CITIES-ADMIN}{战国·秦取韩十城并接管前沿}。十城名称与占领期限没有完整清单，韩国行政网络被切开、赵接上党便会承受更大风险的逻辑却很清楚。前261年王龁攻上党，赵廉颇在长平筑垒，\eventanchor{WQ-0261-WANGHE-SHANGDANG}{战国·王龁攻上党、赵守长平}；秦又以王龁、白起先后统帅调整部署，\eventanchor{WQ-0261-QIN-CHANGPING-COMMAND}{战国·秦调整长平统帅}。主将交接、军令保密程度和每段防线都存异文，双方为粮道而相持则是无法绕开的背景。

前260年昭襄王密令白起取代王龁，\eventanchor{WQ-0260-BAIQI-ASSUMES-COMMAND}{战国·白起接掌长平前线}，秦军转向分割围困赵军。它承接既有的山口、运输、指挥与赵廷判断，不应把战果归成一纸密令的魔力。战后秦廷是否立即强攻邯郸，白起、范雎与昭王之间的分歧，\eventanchor{WQ-0259-QIN-HANDAN-PEACE-DEBATE}{战国·长平后秦廷议邯郸}，人物责任被传记写得格外鲜明，实际争论仍只能审慎复原。白起终因战议失和、被赐死，\eventanchor{WQ-0259-BAIQI-DISMISSAL}{战国·白起失势被赐死}；这不是秦突然失去所有优势，而是连年战争也会使最有声望的将领与王廷对于成本、时机和责任产生不能调和的裂缝。

救赵联军逼近后，秦军由邯郸外线退向河东、上党，\eventanchor{WQ-0257-QIN-HANDAN-RETREAT-ROUTE}{战国·秦军自邯郸外线撤退}。撤军路线和损失有争议，赵能保国、秦仍保有长平后大部前沿的双重结果不应被抹平。前256年前后，李冰主持都江堰的传统年代，\eventanchor{WQ-0256-DUJIANGYAN}{战国·都江堰传统始建年代}，可提示巴蜀水利与关中战略的长远关联；工程分期与主持者仍有争论，不能把今日可见的工程全貌压缩为昭襄王一年的政令。前251年昭襄王卒、孝文王继位，\eventanchor{WQ-0251-ZHAOXIANG-DEATH}{战国·昭襄王卒}。他留下的不是已然统一的天下，而是一套把边区、河道、前沿城邑和朝廷人事反复接合的扩张次序。

\begin{sourcenote}[从人物传记到战争地理]
\sourcebook{资治通鉴}为这一段提供了可连贯对读的编年脊柱；《史记·秦本纪》《白起王翦列传》保存伊阙及秦军推进的基本线索，《范雎蔡泽列传》《穰侯列传》保存范雎与魏冉更替的中心叙事，《赵世家》《韩世家》《魏世家》则从近邻的处境叙述上党与城邑压力。《匈奴列传》《河渠书》《华阳国志》及蜀道、河道、都江堰与战国历史地理研究，分别可检验义渠、巴蜀、水利与道路的空间条件；它们仍不能把路线、建置和工程阶段说得过满。现代研究常以伊阙、野王、上党与太行通道的地理，检验“远交近攻”在粮运和占领上的可行性；同时也提醒我们，范雎游说的完整文本、白起战果的巨数和上党决策的动机，都经过秦汉以后的文本传统塑形。
\end{sourcenote}

\begin{turningpoint}[制度得失：胜利必须接上道路]
昭襄王时期的秦廷把东进的军政资源逐步收束到王权能够主导的次序中，缓解了相权、外戚与军功集团各自牵动外交的风险；代价是近邻国家和边地社会承受了更连续的征发、筑城与战争压力。伊阙打开的不是一条自动通往统一的直线，范雎所概括的也不是无需代价的公式。秦真正获得的，是把一块近处土地、一道山口和下一批军粮接进既有网络的能力；上党危机说明，这种能力一旦同赵的资源和判断相撞，胜利同样可能被拖进一场谁也难以轻易收束的持久战。
\end{turningpoint}
